% 04-monitor-control.tex - Monitor and Control Interface

\chapter{Monitor and Control Interface}
\label{sec:monitor-control}

\section{Overview}

Control and monitoring of the DR is performed by two classes of messages:

\begin{enumerate}
    \item \textbf{Monitoring messages} request system status information from the
          DR via the Management Information Base (MIB).
    \item \textbf{Command messages} request that the DR execute some action.
\end{enumerate}

All messages are exchanged over UDP using the format defined in the MCS Common
ICD (LWA Engineering Memo MCS0005).  If the DR cannot comply with a request, a rejection
response containing an error message (see Appendix~\ref{app:errors}) is
returned.

\section{System Configuration}
\label{sec:config}

DROS2 is configured through the \texttt{defaults\_v2.cfg} file, read on
startup or when manually reinitialized via the \cmd{INI} command
(Section~\ref{cmd:ini}).  The configuration parameters are described in
Section~\ref{sec:configuration}.

\begin{newInVersion}{2.4.1}
The configuration file has been renamed from \texttt{defaults.cfg} to
\texttt{defaults\_v2.cfg}.  The \texttt{MessageOutURL} parameter has been
removed; responses are sent to the originating IP address of each request.
\end{newInVersion}

In a multi-DR deployment (Section~\ref{sec:multi-dr}), each DR instance has its
own configuration directory with a separate \texttt{defaults\_v2.cfg} file.

\section{Time Synchronization}

The DR uses \texttt{ntpd} or \texttt{systemd-timesyncd} for time
synchronization.  The DR also synchronizes its internal clock to station time
on each boot-up.  The \cmd{SYN} command (Section~\ref{cmd:syn}) explicitly re-synchronizes time, as does the \cmd{INI} command (Section~\ref{cmd:ini}).

\section{Timing Restrictions}

\begin{itemize}
    \item The DR supports up to 100 commands per second.
    \item Commands that schedule recording must allow at least 5~seconds between
          receipt and start of recording.
    \item Recordings may not be scheduled to begin within 5~seconds of the
          termination of a prior recording session.
\end{itemize}


%% ===========================================================================
%%  MIB
%% ===========================================================================

\section{MIB Structure}
\label{sec:mib}

The Management Information Base is organized as shown in
Table~\ref{tab:mib-summary}.  Each entry is queried using the \cmd{RPT}
command with the entry label as the data field.

\begin{longtable}{l l l p{5.5cm}}
\caption{DR MIB Structure}
\label{tab:mib-summary} \\
\toprule
\textbf{Index} & \textbf{Label} & \textbf{Type} & \textbf{Description} \\
\midrule
\endfirsthead

\toprule
\textbf{Index} & \textbf{Label} & \textbf{Type} & \textbf{Description} \\
\midrule
\endhead

\bottomrule
\endfoot

\multicolumn{4}{l}{\textbf{1.\ MCS-Required Entries}} \\
1.1     & \mib{SUMMARY}    & Status & Current system status \\
1.2     & \mib{INFO}       & Status & Detailed status information \\
1.3     & \mib{LASTLOG}    & Status & Last log entry \\
1.4     & \mib{SUBSYSTEM}  & Status & Subsystem identifier \\
1.5     & \mib{SERIALNO}   & Status & Serial number \\
1.6     & \mib{VERSION}    & Status & Software version \\
\midrule

\multicolumn{4}{l}{\textbf{2.\ Current Operation}} \\
2.1     & \mib{OP-TYPE}          & Status & Current operation type \\
2.2     & \mib{OP-SCHEDULE}      & Status & Operation start and stop times \\
2.2.1   & \mib{OP-START}         & Status & Operation start time \\
2.2.2   & \mib{OP-STOP}          & Status & Operation expected stop time \\
2.3     & \mib{OP-REFERENCE}     & Status & Command reference number \\
2.4.1   & \mib{OP-TAG}           & Status & Internal file tag \\
2.4.2   & \mib{OP-FORMAT}        & Status & Data format in use \\
2.4.3   & \mib{OP-POSITION}  & Status & File position info \\
2.5.1   & \mib{OP-FILENAME}      & Status & External file name and device \\
2.5.2   & \mib{OP-FILEINDEX}     & Status & File series index \\
\midrule

\multicolumn{4}{l}{\textbf{3.\ Schedule}} \\
3.1     & \mib{SCHEDULE-COUNT}   & Status & Count of scheduled operations \\
3.2.\{n\}   & \mib{SCHEDULE-ENTRY-\{n\}} & Status & \{n\}th scheduled operation \\
\midrule

\multicolumn{4}{l}{\textbf{4.\ Directory}} \\
4.1     & \mib{DIRECTORY-COUNT}  & Status & Count of stored recordings \\
4.2.\{n\}   & \mib{DIRECTORY-ENTRY-\{n\}}& Status & \{n\}th recording entry \\
\midrule

\multicolumn{4}{l}{\textbf{5.\ Storage}} \\
5.1     & \mib{TOTAL-STORAGE}    & Status & Total internal storage (bytes) \\
5.2     & \mib{REMAINING-STORAGE}& Status & Available internal storage (bytes) \\
5.3     & \mib{CONTIGUOUS-STORAGE}& Status & Deprecated; always returns~0 \\
\midrule

\multicolumn{4}{l}{\textbf{6.\ Removable Devices}} \\
6.1     & \mib{DEVICE-COUNT}     & Status & Number of external devices \\
6.2.\{n\}   & \mib{DEVICE-ID-\{n\}}     & Status & Device ID of \{n\}th device \\
6.3.\{n\}   & \mib{DEVICE-STORAGE-\{n\}} & Status & Free space on \{n\}th device \\
6.4.\{n\}   & \mib{DEVICE-BARCODE-\{n\}} & Status & Volume name of \{n\}th device \\
\midrule

\multicolumn{4}{l}{\textbf{7.\ CPU Info}} \\
7.1     & \mib{CPU-COUNT}        & Status & Number of CPU cores \\
7.2.\{n\}   & \mib{CPU-TEMP-\{n\}}       & Status & Temperature of core \{n\} \\
\midrule

\multicolumn{4}{l}{\textbf{8.\ HDD Info}} \\
8.1     & \mib{HDD-COUNT}        & Status & Number of hard drives \\
8.2.\{n\}   & \mib{HDD-TEMP-\{n\}}       & Status & Temperature of HDD \{n\} \\
\midrule

\multicolumn{4}{l}{\textbf{9.\ Data Formats}} \\
9.1     & \mib{FORMAT-COUNT}     & Status & Number of supported formats \\
9.2.\{n\}   & \mib{FORMAT-NAME-\{n\}}    & Status & Name of \{n\}th format \\
9.3.\{n\}   & \mib{FORMAT-PAYLOAD-\{n\}} & Status & UDP payload size of \{n\}th format \\
9.4.\{n\}   & \mib{FORMAT-RATE-\{n\}}    & Status & Data rate of \{n\}th format \\
\midrule

\multicolumn{4}{l}{\textbf{10.\ System Log}} \\
10.1    & \mib{LOG-COUNT}        & Status & Number of log entries \\
10.2.\{n\}  & \mib{LOG-ENTRY-\{n\}}      & Status & \{n\}th log entry \\
\midrule

\multicolumn{4}{l}{\textbf{11.\ Live Buffer}} \\
11.1    & \mib{BUFFER}           & Data   & Deprecated; see Section~\ref{mib:buffer} \\
11.2    & \mib{BUFFER-RESTRICT}  & Data   & Deprecated; see Section~\ref{mib:buffer-restrict} \\
\midrule

\multicolumn{4}{l}{\textbf{12.\ DRSU Status}} \\
12.1    & \mib{DRSU-COUNT}       & Status & Number of detected DRSUs \\
12.2    & \mib{DRSU-SELECTED}    & Status & Currently selected DRSU \\
12.3    & \mib{DRSU-BARCODE}     & Status & Barcode of selected DRSU \\
12.4.\{n\}  & \mib{DRSU-INFO-\{n\}}      & Status & Info on \{n\}th detected DRSU \\
\midrule

\multicolumn{4}{l}{\textbf{13.\ Timing}} \\
13.1    & \mib{TT-LAG}           & Status & Current time-tag lag \\
13.2    & \mib{TT-LAG-INITIAL}   & Status & Initial time-tag lag \\

\end{longtable}


%% ===========================================================================
%%  MIB ENTRIES IN DETAIL
%% ===========================================================================

\section{MIB Entries in Detail}
\label{sec:mib-detail}

\subsection{SUMMARY}
\label{mib:summary}
\begin{description}[style=nextline]
    \item[Index:] 1.1
    \item[Description:]
        Reports the current system status.
    \item[Response Format:] \texttt{<Status>}
    \item[Response Element:]
        Status (ASCII-7-A): One of the values in Section~\ref{sec:states}.
\end{description}

\subsection{INFO}
\label{mib:info}
\begin{description}[style=nextline]
    \item[Index:] 1.2
    \item[Description:]
        Reports detailed status information.  The content varies depending on
        the current system state.
    \item[Response Format:] \texttt{<Info>}
    \item[Response Element:]
        Info (ASCII-256-A): Human-readable status detail.
\end{description}

\subsection{LASTLOG}
\label{mib:lastlog}
\begin{description}[style=nextline]
    \item[Index:] 1.3
    \item[Description:]
        Reports the most recent system log entry.
    \item[Response Format:] \texttt{<Log Entry>}
    \item[Response Element:]
        Log Entry (ASCII-256-A): Last log message.
\end{description}

\subsection{SUBSYSTEM}
\label{mib:subsystem}
\begin{description}[style=nextline]
    \item[Index:] 1.4
    \item[Description:]
        Reports the subsystem reference designator.
    \item[Response Format:] \texttt{<Designator>}
    \item[Response Element:]
        Designator (ASCII-3-A): Subsystem identifier (e.g., \texttt{DR1}).
\end{description}

\subsection{SERIALNO}
\label{mib:serialno}
\begin{description}[style=nextline]
    \item[Index:] 1.5
    \item[Description:]
        Reports the DR serial number as configured in \texttt{defaults\_v2.cfg}.
    \item[Response Format:] \texttt{<Serial Number>}
    \item[Response Element:]
        Serial Number (ASCII-20-A).
\end{description}

\subsection{VERSION}
\label{mib:version}
\begin{description}[style=nextline]
    \item[Index:] 1.6
    \item[Description:]
        Reports the DROS2 software version.
    \item[Response Format:] \texttt{<Version>}
    \item[Response Element:]
        Version (ASCII-256-A): Software version string.
\end{description}

\subsection{OP-TYPE}
\label{mib:op-type}
\begin{description}[style=nextline]
    \item[Index:] 2.1
    \item[Description:]
        Reports the current operation type.  If no operation is in progress,
        reports idle state.
    \item[Response Format:] \texttt{<Operation Type>}
    \item[Response Element:]
        Operation Type (ASCII-11-A) --- one of the values in
        Table~\ref{tab:op-types}.
\end{description}

\begin{table}[h]
\centering
\caption{Operation Type Values}
\label{tab:op-types}
\begin{tabular}{l l}
\toprule
\textbf{Value} & \textbf{Meaning} \\
\midrule
\texttt{Idle}        & No operation in progress \\
\texttt{Record}      & Recording NDP data \\
\texttt{Spectrometr} & Computing spectrometer output \\
\texttt{Copy}        & Copying recording to external storage \\
\texttt{Dump}        & Dumping blocks to external storage \\
\texttt{Format}      & Formatting storage device \\
\bottomrule
\end{tabular}
\end{table}

\subsection{OP-SCHEDULE}
\label{mib:op-schedule}
\begin{description}[style=nextline]
    \item[Index:] 2.2
    \item[Description:]
        Reports the start and stop times of the current operation as a
        single combined response.  Not valid if operation is \texttt{Idle}.
    \item[Response Format:] \texttt{<Start MJD>'<Start MPM>'<Stop MJD>'<Stop MPM>}
    \item[Response Elements:]~\\
        Start MJD (ASCII-6-\#): MJD when operation began. \\
        Start MPM (ASCII-9-\#): MPM when operation began. \\
        Stop MJD (ASCII-6-\#): MJD when operation will end. \\
        Stop MPM (ASCII-9-\#): MPM when operation will end.
\end{description}

\subsection{OP-START}
\label{mib:op-start}
\begin{description}[style=nextline]
    \item[Index:] 2.2.1
    \item[Description:]
        Reports the time at which the current operation began.  Not valid if
        operation is \texttt{Idle}.
    \item[Response Format:] \texttt{<Start MJD>'<Start MPM>}
    \item[Response Elements:]~\\
        Start MJD (ASCII-6-\#): MJD when operation began. \\
        Start MPM (ASCII-9-\#): MPM when operation began.
\end{description}

\subsection{OP-STOP}
\label{mib:op-stop}
\begin{description}[style=nextline]
    \item[Index:] 2.2.2
    \item[Description:]
        Reports the scheduled/expected end time of the current operation.
        Only valid for \texttt{Record}, \texttt{Spectrometr}, \texttt{Copy},
        or \texttt{Dump} operations.
    \item[Response Format:] \texttt{<Stop MJD>'<Stop MPM>}
    \item[Response Elements:]~\\
        Stop MJD (ASCII-6-\#): MJD when operation will end. \\
        Stop MPM (ASCII-9-\#): MPM when operation will end.
\end{description}

\subsection{OP-REFERENCE}
\label{mib:op-reference}
\begin{description}[style=nextline]
    \item[Index:] 2.3
    \item[Description:]
        Reports the reference number of the command message that
        initiated/scheduled the current operation.  Not valid if operation
        is \texttt{Idle}.
    \item[Response Format:] \texttt{<Reference Number>}
    \item[Response Element:]
        Reference Number (ASCII-9-\#).
\end{description}

\subsection{OP-TAG}
\label{mib:op-tag}
\begin{description}[style=nextline]
    \item[Index:] 2.4.1
    \item[Description:]
        Reports the tag value identifying the file in use.  Only valid for
        \texttt{Record}, \texttt{Spectrometr}, \texttt{Copy}, or \texttt{Dump}
        operations.
    \item[Response Format:] \texttt{<Tag>}
    \item[Response Element:]
        Tag (ASCII-16-A): filename tag in the form
        \texttt{<MJD>\_<Reference Number>} including the literal underscore.
\end{description}

\subsection{OP-FORMAT}
\label{mib:op-format}
\begin{description}[style=nextline]
    \item[Index:] 2.4.2
    \item[Description:]
        Reports the data format in use.  For \texttt{Record} operations,
        this is the format being recorded.  For \texttt{Copy}/\texttt{Dump},
        it is the format specified by the originating \cmd{REC} command.
    \item[Response Format:] \texttt{<Data Format>}
    \item[Response Element:]
        Data Format (ASCII-32-A).
\end{description}

\subsection{OP-POSITION}
\label{mib:op-fileposition}
\begin{description}[style=nextline]
    \item[Index:] 2.4.3
    \item[Description:]
        Reports start position, length, and current position of the file in
        use.  Current Position is the offset relative to Start Position.
        Valid for \texttt{Record}, \texttt{Copy}, \texttt{Dump} operations.
        In \texttt{Spectrometr} mode, Start Position and Length are not valid,
        but Current Position reflects the current output file size.
    \item[Response Format:]
        \texttt{<Start Position>'<Length>'<Current Position>}
    \item[Response Elements:]~\\
        Start Position (ASCII-15-\#): Position of first byte (always 0 for Recording). \\
        Length (ASCII-15-\#): Copy: bytes to copy; Dump: chunk size; Record: expected file size. \\
        Current Position (ASCII-15-\#): Position of most recent byte written/read.
\end{description}

\subsection{OP-FILENAME}
\label{mib:op-filename}
\begin{description}[style=nextline]
    \item[Index:] 2.5.1
    \item[Description:]
        Reports the file name and external storage device ID.  Only valid for
        \texttt{Copy} or \texttt{Dump} operations.  For \texttt{Dump}, the
        returned file name is the series name; individual files are named
        \texttt{<Filename>.X} with a zero-padded serial identifier.
    \item[Response Format:] \texttt{<Storage ID>'<Filename>}
    \item[Response Elements:]~\\
        Storage ID (ASCII-64-A): Linux partition or mount point. \\
        Filename (ASCII-128-A): File or series name.
\end{description}

\subsection{OP-FILEINDEX}
\label{mib:op-fileindex}
\begin{description}[style=nextline]
    \item[Index:] 2.5.2
    \item[Description:]
        Reports which file of a file series is being written to.
        Only valid for \texttt{Dump} operations.
    \item[Response Format:] \texttt{<File Index>}
    \item[Response Element:]
        File Index (ASCII-9-\#).
\end{description}

\subsection{SCHEDULE-COUNT}
\label{mib:schedule-count}
\begin{description}[style=nextline]
    \item[Index:] 3.1
    \item[Description:]
        Reports the count of all scheduled operations, including the current
        operation if one is in progress.
    \item[Response Format:] \texttt{<Count>}
    \item[Response Element:]
        Count (ASCII-6-\#).
\end{description}

\subsection{SCHEDULE-ENTRY-\{n\}}
\label{mib:schedule-entry}
\begin{description}[style=nextline]
    \item[Index:] 3.2.\{n\}
    \item[Description:]
        Reports relevant information for the \{n\}th scheduled operation.
    \item[Response Format:]
        \texttt{<Operation Type>'<Reference Number>'<Start MJD>'<Start MPM>'<Stop MJD>'<Stop MPM>'<Data Format>}
    \item[Response Elements:]~\\
        Operation Type (ASCII-11-A). \\
        Reference Number (ASCII-9-\#). \\
        Start MJD (ASCII-6-\#). \\
        Start MPM (ASCII-9-\#). \\
        Stop MJD (ASCII-6-\#). \\
        Stop MPM (ASCII-9-\#). \\
        Data Format (ASCII-32-A).
\end{description}

\subsection{DIRECTORY-COUNT}
\label{mib:directory-count}
\begin{description}[style=nextline]
    \item[Index:] 4.1
    \item[Description:]
        Reports the number of recordings contained on internal storage.
    \item[Response Format:] \texttt{<Count>}
    \item[Response Element:]
        Count (ASCII-6-\#).
\end{description}

\subsection{DIRECTORY-ENTRY-\{n\}}
\label{mib:directory-entry}
\begin{description}[style=nextline]
    \item[Index:] 4.2.\{n\}
    \item[Description:]
        Reports pertinent information for the \{n\}th recording on internal storage.
    \item[Response Format:]
        \texttt{<Tag>'<Start MJD>'<Start MPM>'<Stop MJD>'<Stop MPM>'<Data Format>'<Size>'<Disk Usage>'<Complete>}
    \item[Response Elements:]~\\
        Tag (ASCII-16-A): Form \texttt{<MJD>\_<Reference Number>}. \\
        Start MJD (ASCII-6-\#). \\
        Start MPM (ASCII-9-\#). \\
        Stop MJD (ASCII-6-\#). \\
        Stop MPM (ASCII-9-\#). \\
        Data Format (ASCII-32-A): Format used when recorded. \\
        Size (ASCII-15-\#): Bytes written to disk. \\
        Disk Usage (ASCII-15-\#): Total bytes occupied on disk. \\
        Complete (ASCII-3-A): \texttt{YES} or \texttt{NO~}.
\end{description}

\subsection{TOTAL-STORAGE}
\label{mib:total-storage}
\begin{description}[style=nextline]
    \item[Index:] 5.1
    \item[Description:]
        Reports total storage capacity of internal storage in bytes.  Returns 0
        when storage is offline or unavailable.
    \item[Response Format:] \texttt{<Size>}
    \item[Response Element:]
        Size (ASCII-15-\#): Total internal storage in bytes.
\end{description}

\subsection{REMAINING-STORAGE}
\label{mib:remaining-storage}
\begin{description}[style=nextline]
    \item[Index:] 5.2
    \item[Description:]
        Reports available bytes on internal storage.
    \item[Response Format:] \texttt{<Available>}
    \item[Response Element:]
        Available (ASCII-15-\#): Unused internal storage in bytes.
\end{description}

\subsection{CONTIGUOUS-STORAGE}
\label{mib:contiguous-storage}
\begin{description}[style=nextline]
    \item[Index:] 5.3
    \item[Description:]
        Deprecated.  Always returns~0.  Retained for backward compatibility
        with older MCS implementations.
    \item[Response Format:] \texttt{<Size>}
    \item[Response Element:]
        Size (ASCII-15-\#): Always~0.
\end{description}

\subsection{DEVICE-COUNT}
\label{mib:device-count}
\begin{description}[style=nextline]
    \item[Index:] 6.1
    \item[Description:]
        Reports the number of available external storage devices.
    \item[Response Format:] \texttt{<Count>}
    \item[Response Element:]
        Count (ASCII-6-\#).
\end{description}

\subsection{DEVICE-ID-\{n\}}
\label{mib:device-id}
\begin{description}[style=nextline]
    \item[Index:] 6.2.\{n\}
    \item[Description:]
        Reports the device ID of the \{n\}th external storage device.
    \item[Response Format:] \texttt{<Storage ID>}
    \item[Response Element:]
        Storage ID (ASCII-64-A): Linux partition (e.g., \texttt{/dev/sdf1}).
\end{description}

\subsection{DEVICE-STORAGE-\{n\}}
\label{mib:device-storage}
\begin{description}[style=nextline]
    \item[Index:] 6.3.\{n\}
    \item[Description:]
        Reports free storage space on the \{n\}th external storage device.
    \item[Response Format:] \texttt{<Available>}
    \item[Response Element:]
        Available (ASCII-15-\#): Unused bytes.  Returns 0 if the device is
        not properly formatted or uses an unsupported file system.
\end{description}

\subsection{DEVICE-BARCODE-\{n\}}
\label{mib:device-barcode}
\begin{description}[style=nextline]
    \item[Index:] 6.4.\{n\}
    \item[Description:]
        Reports the volume name (barcode) of the \{n\}th external storage device.
    \item[Response Format:] \texttt{<Volume Name>}
    \item[Response Element:]
        Volume Name (ASCII-64-A): Volume label of the device.
\end{description}

\subsection{CPU-COUNT}
\label{mib:cpu-count}
\begin{description}[style=nextline]
    \item[Index:] 7.1
    \item[Description:]
        Reports the number of CPU cores.
    \item[Response Format:] \texttt{<Count>}
    \item[Response Element:]
        Count (ASCII-3-\#).
\end{description}

\subsection{CPU-TEMP-\{n\}}
\label{mib:cpu-temp}
\begin{description}[style=nextline]
    \item[Index:] 7.2.\{n\}
    \item[Description:]
        Reports the temperature of CPU core \{n\}.
    \item[Response Format:] \texttt{<Core \{n\} Temp>}
    \item[Response Element:]
        Core \{n\} Temp (ASCII-3-\#): Temperature in degrees Celsius.
\end{description}

\subsection{HDD-COUNT}
\label{mib:hdd-count}
\begin{description}[style=nextline]
    \item[Index:] 8.1
    \item[Description:]
        Reports the number of hard drives comprising internal storage.
    \item[Response Format:] \texttt{<Count>}
    \item[Response Element:]
        Count (ASCII-3-\#).
\end{description}

\subsection{HDD-TEMP-\{n\}}
\label{mib:hdd-temp}
\begin{description}[style=nextline]
    \item[Index:] 8.2.\{n\}
    \item[Description:]
        Reports the temperature of the \{n\}th hard drive.  Some hardware may not
        support temperature reporting; in such cases the response may be empty.
    \item[Response Format:] \texttt{<HDD \{n\} Temp>}
    \item[Response Element:]
        HDD \{n\} Temp (ASCII-3-\#): Temperature in degrees Celsius.
\end{description}

\subsection{FORMAT-COUNT}
\label{mib:format-count}
\begin{description}[style=nextline]
    \item[Index:] 9.1
    \item[Description:]
        Returns the number of recording formats supported.
    \item[Response Format:] \texttt{<Count>}
    \item[Response Element:]
        Count (ASCII-6-\#).
\end{description}

\subsection{FORMAT-NAME-\{n\}}
\label{mib:format-name}
\begin{description}[style=nextline]
    \item[Index:] 9.2.\{n\}
    \item[Description:]
        Returns the name of the \{n\}th recording format.
    \item[Response Format:] \texttt{<Format Name>}
    \item[Response Element:]
        Format Name (ASCII-32-A).
\end{description}

\subsection{FORMAT-PAYLOAD-\{n\}}
\label{mib:format-payload}
\begin{description}[style=nextline]
    \item[Index:] 9.3.\{n\}
    \item[Description:]
        Returns the UDP packet payload size of the \{n\}th recording format.
    \item[Response Format:] \texttt{<UDP Payload Size>}
    \item[Response Element:]
        UDP Payload Size (ASCII-4-\#): Bytes of payload.
\end{description}

\subsection{FORMAT-RATE-\{n\}}
\label{mib:format-rate}
\begin{description}[style=nextline]
    \item[Index:] 9.4.\{n\}
    \item[Description:]
        Returns the data rate of the \{n\}th recording format in bytes per second.
    \item[Response Format:] \texttt{<Rate>}
    \item[Response Element:]
        Rate (ASCII-9-\#): Overall data rate (MAC/UDP headers excluded).
\end{description}

\subsection{LOG-COUNT}
\label{mib:log-count}
\begin{description}[style=nextline]
    \item[Index:] 10.1
    \item[Description:]
        Reports the number of system log entries.
    \item[Response Format:] \texttt{<Count>}
    \item[Response Element:]
        Count (ASCII-6-\#).
\end{description}

\subsection{LOG-ENTRY-\{n\}}
\label{mib:log-entry}
\begin{description}[style=nextline]
    \item[Index:] 10.2.\{n\}
    \item[Description:]
        Reports the \{n\}th entry in the system log.
    \item[Response Format:]
        \texttt{<MJD>'<MPM>'<Message Class>'<Message>}
    \item[Response Elements:]~\\
        MJD (ASCII-6-\#): MJD when entry was logged. \\
        MPM (ASCII-9-\#): MPM when entry was logged. \\
        Message Class (ASCII-7-A): \texttt{info~~~}, \texttt{warning}, or
            \texttt{error~~} (trailing spaces pad to 7~characters). \\
        Message (ASCII-234-A): Human-readable description ($\le 234$~characters).
\end{description}

\subsection{BUFFER}
\label{mib:buffer}
\begin{description}[style=nextline]
    \item[Index:] 11.1
    \item[Description:]
        Deprecated.  Returns a deprecation notice.  This entry was previously
        used to return the contents of the live capture buffer.

\begin{calloutBox}{Deprecation Notice}{orange}
The \mib{BUFFER} MIB entry is deprecated in the current DROS2 implementation.
Queries will receive a rejection response.
\end{calloutBox}
\end{description}

\subsection{BUFFER-RESTRICT}
\label{mib:buffer-restrict}
\begin{description}[style=nextline]
    \item[Index:] 11.2
    \item[Description:]
        Deprecated.  Returns a deprecation notice.  This entry was previously
        used to return live capture buffer contents limited to one data packet.

\begin{calloutBox}{Deprecation Notice}{orange}
The \mib{BUFFER-RESTRICT} MIB entry is deprecated in the current DROS2 implementation.
Queries will receive a rejection response.
\end{calloutBox}
\end{description}

\subsection{DRSU-COUNT}
\label{mib:drsu-count}
\begin{description}[style=nextline]
    \item[Index:] 12.1
    \item[Description:]
        Reports the number of detected DRSUs.
    \item[Response Format:] \texttt{<Count>}
    \item[Response Element:]
        Count (ASCII-2-\#).
\end{description}

\subsection{DRSU-SELECTED}
\label{mib:drsu-selected}
\begin{description}[style=nextline]
    \item[Index:] 12.2
    \item[Description:]
        Reports which DRSU is currently active.  If internal storage is down,
        reflects the last valid selection.
    \item[Response Format:] \texttt{<DRSU Number>}
    \item[Response Element:]
        DRSU Number (ASCII-2-\#).
\end{description}

\subsection{DRSU-BARCODE}
\label{mib:drsu-barcode}
\begin{description}[style=nextline]
    \item[Index:] 12.3
    \item[Description:]
        Reports the barcode (volume name) of the currently selected DRSU.
    \item[Response Format:] \texttt{<Barcode>}
    \item[Response Element:]
        Barcode (ASCII-64-A): Volume label of the selected DRSU.
\end{description}

\subsection{DRSU-INFO-\{n\}}
\label{mib:drsu-info}
\begin{description}[style=nextline]
    \item[Index:] 12.4.\{n\}
    \item[Description:]
        Reports information on the \{n\}th detected DRSU.
    \item[Response Format:] \texttt{<Name>'<Partition>'<Unformatted Size>}
    \item[Response Elements:]~\\
        Name (ASCII-6-A): DRSU name (e.g., \texttt{DRSU00}). \\
        Partition (ASCII-64-A): Partition on which the DRSU resides. \\
        Unformatted Size (ASCII-16-\#): Size in bytes before formatting.
\end{description}

\subsection{TT-LAG}
\label{mib:tt-lag}
\begin{description}[style=nextline]
    \item[Index:] 13.1
    \item[Description:]
        Reports the current time-tag lag of the data receiver.  The lag is
        the difference between the expected and observed time tags in the
        incoming NDP data stream, expressed in NDP clock ticks.
    \item[Response Format:] \texttt{<Lag>}
    \item[Response Element:]
        Lag (ASCII-15-\#): Current time-tag lag in clock ticks.
\end{description}

\subsection{TT-LAG-INITIAL}
\label{mib:tt-lag-initial}
\begin{description}[style=nextline]
    \item[Index:] 13.2
    \item[Description:]
        Reports the time-tag lag measured at the start of the most recent
        recording session.
    \item[Response Format:] \texttt{<Lag>}
    \item[Response Element:]
        Lag (ASCII-15-\#): Initial time-tag lag in clock ticks.
\end{description}


%% ===========================================================================
%%  COMMANDS
%% ===========================================================================

\section{Control Commands Summary}
\label{sec:commands}

Table~\ref{tab:commands} lists the control commands supported by the DR.

\begin{longtable}{l l l}
\caption{DR Control Commands}
\label{tab:commands} \\
\toprule
\textbf{Command} & \textbf{Description} & \textbf{Section} \\
\midrule
\endfirsthead

\toprule
\textbf{Command} & \textbf{Description} & \textbf{Section} \\
\midrule
\endhead

\bottomrule
\endfoot

\cmd{INI} & Initialize / restore to boot-up state         & \ref{cmd:ini} \\
\cmd{REC} & Schedule a recording                          & \ref{cmd:rec} \\
\cmd{DEL} & Delete a recording                            & \ref{cmd:del} \\
\cmd{STP} & Stop a scheduled or in-progress operation     & \ref{cmd:stp} \\
\cmd{GET} & Retrieve portion of a recording               & \ref{cmd:get} \\
\cmd{CPY} & Copy recording to external storage            & \ref{cmd:cpy} \\
\cmd{DMP} & Dump recording to file series                 & \ref{cmd:dmp} \\
\cmd{FMT} & Format storage device                         & \ref{cmd:fmt} \\
\cmd{DWN} & Bring internal storage offline                & \ref{cmd:dwn} \\
\cmd{UP}  & Bring internal storage online                 & \ref{cmd:up}  \\
\cmd{SEL} & Select DRSU                                   & \ref{cmd:sel} \\
\cmd{BUF} & Set up live capture buffer                    & \ref{cmd:buf} \\
\cmd{SYN} & Synchronize with NTP server                   & \ref{cmd:syn} \\
\cmd{SPC} & Schedule spectrometer recording               & \ref{cmd:spc} \\
\cmd{SHT} & Shut down DROS2                               & \ref{cmd:sht} \\
\cmd{PNG} & Ping (status query)                           & \ref{cmd:png} \\
\cmd{RPT} & Report (MIB query)                            & \ref{cmd:rpt} \\
\cmd{TST} & Self-test (development only)                  & \ref{cmd:tst} \\

\end{longtable}


%% ===========================================================================
%%  COMMANDS IN DETAIL
%% ===========================================================================

\section{Control Commands in Detail}
\label{sec:commands-detail}

For each command below, the argument format and response format are described.
If no arguments are specified, none are required.  If no response format is
described, the R-COMMENT field is empty on success.  Failed commands return
\texttt{R} in R-RESPONSE with an error message in R-COMMENT
(see Appendix~\ref{app:errors}).

% ---- INI ----
\subsection{INI --- Initialize}
\label{cmd:ini}
\begin{description}[style=nextline]
    \item[Description:]
        Restores the DR to its initial boot-up state except for the system log
        and internal storage contents.
    \item[Argument Format:] \texttt{<Flags>}
    \item[Argument:]
        Flags (ASCII-256-A): Specify \texttt{-{}-flush-log} or \texttt{-L} to
        force system log re-initialization; \texttt{-{}-flush-data} or
        \texttt{-D} to force internal storage re-initialization.  The field
        need not be padded; flag order does not matter.
\end{description}

% ---- REC ----
\subsection{REC --- Record}
\label{cmd:rec}
\begin{description}[style=nextline]
    \item[Description:]
        Schedules or initiates recording of NDP data.  Upon success, returns a
        tag value uniquely identifying the recorded file.

    \item[Argument Format:]
        \texttt{<Start MJD>'<Start MPM>'<Length>'<Data Format>}

    \item[Arguments:]~\\
        Start MJD (ASCII-6-\#): Modified Julian Day; not $>24$~hours in the future. \\
        Start MPM (ASCII-9-\#): Milliseconds past midnight; not within 5~s of
            a prior operation's termination. \\
        Length (ASCII-9-\#): Duration in milliseconds. \\
        Data Format (ASCII-32-A): Pre-configured format name (see
            Section~\ref{mib:format-name} and Appendix~\ref{app:data-formats}).

    \item[Response Format:] \texttt{<Tag>}
    \item[Response Element:]
        Tag (ASCII-16-A): Form \texttt{<MJD>\_<Reference Number>}.
\end{description}

\begin{newInVersion}{2.4.1}
The supported data format names have been updated.  The following formats are
available:
\texttt{DEFAULT\_DRX},
\texttt{DRX\_FILT\_1} through \texttt{DRX\_FILT\_7},
\texttt{DEFAULT\_DRX8},
\texttt{DRX8\_FILT\_1} through \texttt{DRX8\_FILT\_7},
\texttt{DEFAULT\_TBT},
\texttt{DEFAULT\_TBS},
\texttt{TBS\_FILT\_7} through \texttt{TBS\_FILT\_9},
\texttt{DEFAULT\_COR}.
The previously supported \texttt{DEFAULT\_TBN} and \texttt{DEFAULT\_TBW} formats
have been removed.  See Appendix~\ref{app:data-formats} for details.
\end{newInVersion}

\begin{noteBox}
\textbf{Example:} To schedule a 30-second DRX recording starting at MJD~60700,
MPM~43200000 with reference number 42:

\texttt{REC 060700 043200000 000030000 DEFAULT\_DRX}

Response: \texttt{060700\_000000042}
\end{noteBox}

% ---- DEL ----
\subsection{DEL --- Delete}
\label{cmd:del}
\begin{description}[style=nextline]
    \item[Description:]
        Deletes a recording from internal storage.
    \item[Argument Format:] \texttt{<Tag>}
    \item[Argument:]
        Tag (ASCII-16-A): Form \texttt{<MJD>\_<Reference Number>}.
\end{description}

% ---- STP ----
\subsection{STP --- Stop}
\label{cmd:stp}
\begin{description}[style=nextline]
    \item[Description:]
        Halts or prevents a specified operation.  If the operation is scheduled
        but not yet in progress, it is deleted from the schedule and disk space
        is freed.  If it is in progress, it is halted and the file is closed
        but not deleted.  Also stops scheduled or in-progress spectrometer
        operations.
    \item[Argument Format:] \texttt{<Tag>}
    \item[Argument:]
        Tag (ASCII-16-A): Form \texttt{<MJD>\_<Reference Number>}.
\end{description}

% ---- GET ----
\subsection{GET --- Get}
\label{cmd:get}
\begin{description}[style=nextline]
    \item[Description:]
        Retrieves a portion of a specified recording.
    \item[Argument Format:] \texttt{<Tag>'<Start Byte>'<Length>}
    \item[Arguments:]~\\
        Tag (ASCII-16-A): Form \texttt{<MJD>\_<Reference Number>}. \\
        Start Byte (ASCII-15-\#): Byte offset within the file. \\
        Length (ASCII-15-\#): Bytes to return; limited to 8146~bytes.
    \item[Response Format:] \texttt{<Data>}
    \item[Response Element:]
        Data (uint8)$\times$Length: Requested bytes from the file.
\end{description}

% ---- CPY ----
\subsection{CPY --- Copy}
\label{cmd:cpy}
\begin{description}[style=nextline]
    \item[Description:]
        Copies a portion of a recording to a file on external storage.
        Existing files are overwritten without warning.  Not available if
        recordings are scheduled.
    \item[Argument Format:]
        \texttt{<Tag>'<Start Byte>'<Length>'<Device ID>'<Filename>}
    \item[Arguments:]~\\
        Tag (ASCII-16-A): Form \texttt{<MJD>\_<Reference Number>}. \\
        Start Byte (ASCII-15-\#): Byte offset within the file. \\
        Length (ASCII-15-\#): Bytes to copy. \\
        Storage ID (ASCII-64-A): Target device partition. \\
        Filename (ASCII-128-A): Destination file name.
\end{description}

% ---- DMP ----
\subsection{DMP --- Dump}
\label{cmd:dmp}
\begin{description}[style=nextline]
    \item[Description:]
        Copies data blocks from a recording to a series of files on external
        storage.  Files are named \texttt{<Filename>.X} with a zero-padded
        serial identifier.  Not available if recordings are scheduled.
    \item[Argument Format:]
        \texttt{<Tag>'<Start Byte>'<Length>'<Block Size>'<Device ID>'<Filename>}
    \item[Arguments:]~\\
        Tag (ASCII-16-A): Form \texttt{<MJD>\_<Reference Number>}. \\
        Start Byte (ASCII-15-\#): Byte offset within the file. \\
        Length (ASCII-15-\#): Bytes to copy. \\
        Block Size (ASCII-15-\#): Bytes per output file. \\
        Storage ID (ASCII-64-A): Target device partition. \\
        Filename (ASCII-128-A): Series name.
\end{description}

% ---- FMT ----
\subsection{FMT --- Format}
\label{cmd:fmt}
\begin{description}[style=nextline]
    \item[Description:]
        Formats an internal or external storage device.  This is a destructive
        operation that erases all data on the target device.  External
        formatting may require substantial time proportional to device size.
    \item[Argument Format:] \texttt{[Storage ID]}
    \item[Argument:]
        Storage ID (ASCII-64-A): Optional external device partition.
        If omitted, internal storage is formatted.
\end{description}

% ---- DWN ----
\subsection{DWN --- Down}
\label{cmd:dwn}
\begin{description}[style=nextline]
    \item[Description:]
        Prepares internal storage for removal or replacement.  Executes
        immediately; requires a few seconds for completion.  Wait at least
        1~minute after \cmd{DWN} before physically disconnecting storage.
\end{description}

% ---- UP ----
\subsection{UP --- Up}
\label{cmd:up}
\begin{description}[style=nextline]
    \item[Description:]
        Scans for internal storage and brings it online.  If file system
        information cannot be determined, the command is rejected.
\end{description}

% ---- SEL ----
\subsection{SEL --- Select DRSU}
\label{cmd:sel}
\begin{description}[style=nextline]
    \item[Description:]
        Selects an alternate DRSU for internal storage.  The DR must be idle
        with no scheduled operations.  DRSUs must have been previously detected
        by DROS2 (see \mib{DRSU-INFO-\{n\}}).
    \item[Argument Format:] \texttt{<DRSU Identifier>}
    \item[Argument:]
        DRSU Identifier: Either a 0-based numeric index (DRSU to activate;
        0~= first, 1~= second, etc.) or a DRSU barcode string matching the
        volume name reported by \mib{DRSU-BARCODE}.
\end{description}

% ---- BUF ----
\subsection{BUF --- Live Capture Buffer Setup}
\label{cmd:buf}
\begin{description}[style=nextline]
    \item[Description:]
        Initializes the live capture buffer parameters for the next recording
        start.  The buffer stores up to 100~ms of data and can trigger once
        or periodically.  Parameters are reused with subsequent recordings.
        Buffer contents are overwritten unless retrieved via the
        \mib{BUFFER} or \mib{BUFFER-RESTRICT} MIB entries; if retrieved,
        contents are held and subsequent captures go to a shadow buffer.
    \item[Argument Format:]
        \texttt{<OffsetType><WidthType><Retrigger><Order><Behaviour>'<Offset>'<Width>'<Interval>}
    \item[Arguments:]~\\
        OffsetType (ASCII-1-A): \texttt{T} = milliseconds, \texttt{I} = packet
            counts (applies to Offset and Interval). \\
        WidthType (ASCII-1-A): \texttt{T} = milliseconds, \texttt{I} = packet
            counts (applies to Width). \\
        Retrigger (ASCII-1-A): \texttt{A} = periodic, \texttt{N} = one-shot. \\
        Order (ASCII-1-A): Deprecated; supply \texttt{A} or \texttt{B}
            (ignored). \\
        Behaviour (ASCII-1-A): \texttt{D} = delete/reuse, \texttt{H} =
            hold/reuse, \texttt{R} = hold/no-reuse. \\
        Offset (ASCII-16-\#): Milliseconds or packets from recording start. \\
        Width (ASCII-16-\#): Capture length in milliseconds or packets. \\
        Interval (ASCII-16-\#): Milliseconds or packets between captures
            (when Retrigger = \texttt{A}); should be much greater than Width.
\end{description}


\begin{calloutBox}{Implementation Note}{blue}
The \cmd{BUF} command is only functional when DROS2 is built in LiveBuffer mode.
\end{calloutBox}

% ---- SYN ----
\subsection{SYN --- Synchronize}
\label{cmd:syn}
\begin{description}[style=nextline]
    \item[Description:]
        Explicitly synchronizes the DR clock with the station NTP time server.
        Executing while operations are scheduled or in progress may result in
        recording more or less data than desired.
\end{description}

% ---- SPC ----
\subsection{SPC --- Spectrometer}
\label{cmd:spc}

\begin{description}[style=nextline]
    \item[Description:]
        Schedules an integrated spectra recording session.  The spectrometer
        computes power spectra for each tuning and polarization of beam data
        (NDP DRX or DRX8 output) without windowing.  Non-beam data formats
        (TBT, TBS, COR) are not supported by the spectrometer.  Output is stored on the internal
        storage under the spectrometer directory.  The tag value uniquely
        identifies the output file.  The file format is defined in
        Appendix~\ref{app:spectrometer}.

    \item[Argument Format:]
        \texttt{<Start MJD>'<Start MPM>'<Length>'<Transform Length>'<Integration Count>'[Minimum Fill]'[Highwater Threshold]}

    \item[Arguments:]~\\
        Start MJD (ASCII-6-\#): Modified Julian Day; not $>24$~hours in the future. \\
        Start MPM (ASCII-9-\#): Milliseconds past midnight. \\
        Length (ASCII-9-\#): Duration in milliseconds.  Same restrictions as
            \cmd{REC}. \\
        Transform Length (ASCII-6-\#): FFT length in channels.
            Valid range: 32--16384. \\
        Integration Count (ASCII-8-\#): Number of consecutive spectra to
            integrate.  Valid range: 384--24576. \\
        Minimum Fill (ASCII-8-\#): Optional (default = Integration Count / 2).
            Minimum number of samples in a ``cell'' for valid output.
            Range: 2 to Integration Count. \\
        Highwater Threshold (ASCII-4-\#): Optional (default = 128).
            Minimum buffered ``blocks'' before partial processing.

    \item[Validation:]
        The product (Transform Length $\times$ Integration Count) must be an
        integral multiple of 4096 (the DRX samples-per-frame count).

    \item[Response Format:] \texttt{<Tag>}
    \item[Response Element:]
        Tag (ASCII-16-A): Form \texttt{<MJD>\_<Reference Number>}.

\end{description}

\begin{calloutBox}{Implementation Note}{blue}
The \cmd{SPC} command is only available when DROS2 is built in Spectrometer mode.
\end{calloutBox}

\begin{newInVersion}{2.4.1}
The spectrometer has been significantly enhanced since version~1.4:
\begin{itemize}
    \item \textbf{Stokes products:} The output Stokes product is selected via
        the \texttt{Stokes} metadata field in the MCS message.
        Supported values: \texttt{XXYY} (default), \texttt{CRCI},
        \texttt{IQUV}, \texttt{IV}, \texttt{XXCRCIYY}, \texttt{I}.
    \item \textbf{Flexible FFT parameters:} The Transform Length and Integration
        Count are now fully configurable within the supported ranges
        (32--16384 channels, 384--24576 integrations).  The previous restriction
        to 32~channels and 6144~integrations has been removed.
    \item \textbf{DRX and DRX8 input:} The spectrometer accepts both DRX
        (4-bit) and DRX8 (8-bit) input data.
\end{itemize}
\end{newInVersion}

\begin{noteBox}
\textbf{Example:} To schedule a 60-second spectrometer session at MJD~60700,
MPM~43200000 with 256~channels, 4096~integrations, and XXYY Stokes output:

\texttt{SPC 060700 043200000 000060000 000256 00004096}

Response: \texttt{060700\_000000042}
\end{noteBox}

% ---- SHT ----
\subsection{SHT --- Shutdown}
\label{cmd:sht}
\begin{description}[style=nextline]
    \item[Description:]
        Initiates a shutdown of DROS2.  The system state transitions
        to \texttt{SHUTDWN} and all resources are released.  When deployed as
        a \texttt{systemd} service, the process will be restarted automatically
        after the configured delay.
    \item[Argument Format:] \texttt{[Mode]}
    \item[Argument:]
        Mode (ASCII-7-A): Optional.  \texttt{SCRAM} for an immediate shutdown
        that bypasses normal cleanup.  \texttt{RESTART} to request a clean
        restart.  If omitted, an orderly shutdown is performed.
\end{description}

% ---- PNG ----
\subsection{PNG --- Ping}
\label{cmd:png}
\begin{description}[style=nextline]
    \item[Description:]
        Returns the current system state in the R-COMMENT field.  Used as a
        simple connectivity and status check.
    \item[Response Format:] \texttt{<System State>}
    \item[Response Element:]
        System State (ASCII-7-A): One of the values in
        Section~\ref{sec:states}.
\end{description}

% ---- RPT ----
\subsection{RPT --- Report}
\label{cmd:rpt}
\begin{description}[style=nextline]
    \item[Description:]
        Queries a MIB entry and returns its value.  The MIB entry label is
        supplied as the data field of the command message.
    \item[Argument Format:] \texttt{<MIB Label>}
    \item[Argument:]
        MIB Label: The label of the MIB entry to query (see
        Section~\ref{sec:mib}).
    \item[Response:]
        The response data field contains the MIB entry value as described in
        Section~\ref{sec:mib-detail}.
\end{description}

% ---- TST ----
\subsection{TST --- Self Test}
\label{cmd:tst}
\begin{description}[style=nextline]
    \item[Description:]
        Performs a system self-test.  This command is for development purposes
        only and is not supported by this ICD.
\end{description}


%% ===========================================================================
%%  SYSTEM STATES
%% ===========================================================================

\section{System States}
\label{sec:states}

The DR reports one of the following system states in response to \cmd{PNG}
commands and in the R-SUMMARY field of all responses:

\begin{longtable}{l p{10cm}}
\toprule
\textbf{State} & \textbf{Description} \\
\midrule
\endfirsthead
\endhead
\bottomrule
\endfoot

\texttt{SHUTDWN} &
    The DR is in the process of shutting down. \\

\texttt{BOOTING} &
    The DR is initializing (scanning storage, starting services). \\

\texttt{NORMAL~} &
    The DR is operating normally and ready to accept commands. \\

\texttt{WARNING} &
    The DR has detected a non-critical condition (e.g., elevated temperature,
    low storage space). \\

\texttt{ERROR~~} &
    The DR has detected a critical error preventing normal operation. \\

\end{longtable}
